\section{Discussão} \label{sec:discussao}

\subsection{Achados Principais, Estudo-Base e Recomendações}

Nos estudos avaliados, há contraste claro entre modelos avançados e rigor:
\nFamDL/\nTotal\ (\pctFamDL) usam \textit{deep learning}, mas só
\nRiscoBaixo/\nAplicaveisVaz\ (\pctRiscoBaixoAplic) mantêm risco baixo de
vazamento. O artigo-base~\cite{varghese_machine_2024} adotou CV aninhada,
acurácia balanceada, código público e PD$\times$DD, mas poucos mantiveram esses
critérios (\nValAninhada/\nTotal, \nMetBalAcc/\nTotal\ e
\nCodDisponivel/\nTotal, nessa ordem). Protocolos sem controle ou com
balanceamento pré-\textit{split}~%
\cite{abioye_improving_2025,wang_parkinsons_2025,yu_improved_2024} chegam a
92\%--100\%, incomparáveis à divisão por
sujeito~\cite{zeebaree_three-class_2027,vu_simple_2026}, conforme alertam
revisões~\cite{pinto_neto_machine_2026,ouabd_overview_2026}.

Para avançar a área, quatro cuidados são prioritários: separar por participante
antes de janelar, normalizar ou gerar amostras; usar métricas balanceadas
(Bal.\ Acc., F1-macro, MCC) e calibração; focar tarefas com DD; e abrir código,
lendo \textit{preprints} (\nPreprints) com cautela.

\subsection{Ameaças à Validade}

Na busca (I7), focar citações ao artigo-base~\cite{varghese_machine_2024} pode
omitir quem cita só o
PhysioNet~\cite{PhysioNet-parkinsons-disease-smartwatch-1.0.0}, além de faltarem
corte e rejeições. Na extração (I2), IA sem dupla revisão nem $\kappa$ deixa
\nPontosRevisao\ pontos e metadados pendentes. Os \textit{preprints} (B4) foram
tratados na Seção~\ref{sec:preprints}; QA e vazamento (I4, I5) avaliam relato, e
a validação externa (I1) requer revisão em \nMultiBase/\nTotal\ multibase
e~\cite{abioye_improving_2025}.
